\clearpage % Bắt buộc để ngắt hẳn sang trang mới, không bị dính vào trang trước
\phantomsection
\addcontentsline{toc}{chapter}{THÔNG TIN KẾT QUẢ NGHIÊN CỨU ĐỀ TÀI}

\noindent
\begin{tabular}{@{}c@{}}
    \fontsize{12pt}{15pt}\selectfont BỘ GIÁO DỤC VÀ ĐÀO TẠO \\
    \fontsize{12pt}{15pt}\selectfont \textbf{ĐẠI HỌC CẦN THƠ} \\[-20pt] % Kéo đường kẻ sát lên dòng trên
    \rule{3.2cm}{0.6pt}
\end{tabular}
\vspace{0.8cm}

% --- TIÊU ĐỀ TRANG CĂN GIỮA ---
\begin{center}
    {\fontsize{14pt}{18pt}\selectfont\bfseries THÔNG TIN KẾT QUẢ NGHIÊN CỨU ĐỀ TÀI\par}
\end{center}

\vspace{0.5cm}

% --- NỘI DUNG CHI TIẾT ---
\noindent\textbf{1. Thông tin chung:}
\begin{itemize}[leftmargin=0.8cm, label=\textbullet, itemsep=2pt, topsep=2pt]
    \item Mã số đề tài: TSV2026-127.
    \item Tên đề tài: Hệ thống nhận diện thực phẩm và xác định giá trị dinh dưỡng ứng dụng Deep Learning.
    \item Sinh viên chủ nhiệm đề tài: Nguyễn Hữu Thanh Hòa; lớp DI23Z6A1; ngành Khoa học Máy tính.
    \item Người hướng dẫn: PGS. TS. Nguyễn Thái Nghe.
\end{itemize}

\vspace{0.3cm}
\noindent\textbf{2. Mục tiêu đề tài:}
\begin{itemize}[leftmargin=0.8cm, label=\textbullet, itemsep=2pt, topsep=2pt]
    \item Xây dựng hệ thống thông minh ứng dụng học sâu (Deep Learning) có khả năng nhận diện chính xác các loại thực phẩm từ ảnh chụp đơn lẻ.
    \item Xây dựng cơ sở dữ liệu dinh dưỡng và y tế tích hợp, kết hợp mô hình Lightweight OCR đọc khối lượng trên khay thực phẩm để tính toán tổng giá trị dinh dưỡng của bữa ăn.
    \item Phát triển thuật toán phân tích, cảnh báo mức độ an toàn và khuyến nghị món ăn phù hợp với bệnh nhân béo phì, cao huyết áp trên ứng dụng di động NutriCare.
\end{itemize}

\vspace{0.3cm}
\noindent\textbf{3. Tính mới và sáng tạo:}
\begin{itemize}[leftmargin=0.8cm, label=\textbullet, itemsep=2pt, topsep=2pt]
    \item Kết hợp mô hình Deep Learning phân loại thực phẩm với mô hình YOLO định vị và mô hình OCR nhẹ nhằm tự động hóa việc xác định khối lượng từ ảnh chụp nhãn tem thực tế.
    \item Xây dựng cơ sở dữ liệu dinh dưỡng đặc thù cho thực phẩm Việt Nam gắn liền với bảng thành phần của Viện Dinh dưỡng Quốc gia.
    \item Tích hợp hệ thống luật logic y tế tại Backend để tự động đối chiếu chỉ số vi chất và đưa ra cảnh báo bệnh lý theo thời gian thực thay vì người dùng phải nhập thủ công.
\end{itemize}

\vspace{0.3cm}
\noindent\textbf{4. Kết quả nghiên cứu:}
\begin{itemize}[leftmargin=0.8cm, label=\textbullet, itemsep=2pt, topsep=2pt]
    \item Xây dựng và chuẩn hóa tập dữ liệu hình ảnh thực phẩm Việt Nam kèm nhãn dinh dưỡng và tập ảnh tem nhãn khay phục vụ nhận dạng ký tự.
    \item Huấn luyện mô hình YOLO định vị nhãn và mô hình học sâu phân loại thực phẩm đạt độ chính xác cao, tối ưu tốc độ xử lý trên thiết bị di động.
    \item Hoàn thiện dịch vụ Backend API xử lý đồng thời luồng nhận diện món ăn, bóc tách OCR khối lượng, tra cứu vi chất và kiểm tra quy tắc y tế.
    \item Phát triển thành công nguyên mẫu ứng dụng di động NutriCare cho phép thiết lập hồ sơ bệnh lý, chụp ảnh món ăn và hiển thị bảng phân tích dinh dưỡng trực quan.
\end{itemize}

\vspace{0.3cm}
\noindent\textbf{5. Sản phẩm:}
\begin{itemize}[leftmargin=0.8cm, label=\textbullet, itemsep=2pt, topsep=2pt]
    \item Ứng dụng di động NutriCare hoàn chỉnh trên nền tảng Android/iOS.
    \item Tệp trọng số mô hình Deep Learning nhận diện thực phẩm, mô hình YOLO và Lightweight OCR.
    \item Mã nguồn hệ thống Backend API và cơ sở dữ liệu dinh dưỡng y học tích hợp.
    \item Báo cáo tổng kết đề tài, báo cáo tóm tắt, bản tin và video clip giới thiệu đề tài.
\end{itemize}

\vspace{0.3cm}
\noindent\textbf{6. Công bố khoa học:}
\begin{itemize}[leftmargin=0.8cm, label=\textbullet, itemsep=2pt, topsep=2pt]
    \item[] Tại thời điểm hoàn thiện báo cáo, đề tài chưa ghi nhận công bố khoa học đã được chấp nhận.
\end{itemize}

\vspace{0.3cm}
\noindent\textbf{7. Đóng góp và khả năng áp dụng:}
\begin{itemize}[leftmargin=0.8cm, label=\textbullet, itemsep=2pt, topsep=2pt]
    \item Đóng góp giải pháp kết hợp thị giác máy tính và OCR để giải quyết bài toán ước lượng khẩu phần ăn tự động phù hợp với điều kiện Việt Nam.
    \item Hỗ trợ người dùng cá nhân (nhất là người mắc bệnh béo phì, cao huyết áp) chủ động kiểm soát năng lượng nạp vào và duy trì lối sống lành mạnh.
    \item Có thể mở rộng triển khai tại các phòng khám dinh dưỡng, bệnh viện hoặc tích hợp vào các nền tảng y tế số (HealthTech).
\end{itemize}

\vspace{0.3cm}
\noindent\textbf{8. Hiệu quả và phương thức chuyển giao:}
\begin{itemize}[leftmargin=0.8cm, label=\textbullet, itemsep=2pt, topsep=2pt]
    \item Tiết kiệm thời gian ghi chép và hạn chế sai số so với phương pháp nhập liệu bữa ăn thủ công trên các ứng dụng truyền thống.
    \item Chuyển giao dưới dạng gói mã nguồn mô-đun hóa, tệp cài đặt ứng dụng di động và tài liệu đặc tả API hệ thống.
\end{itemize}

\vspace{0.6cm}
\noindent\textbf{Nhận xét của người hướng dẫn về những đóng góp khoa học của sinh viên thực hiện đề tài:}\par
\vspace{0.2cm}
\noindent\dotfill\par
\vspace{0.1cm}
\noindent\dotfill\par
\vspace{0.1cm}
\noindent\dotfill\par
\vspace{0.1cm}
\noindent\dotfill\par

\vspace{0.8cm}
\hfill
\begin{tabular}{c}
    \textit{Cần Thơ, ngày \dots\dots\ tháng 10 năm 2026} \\
    \textbf{NGƯỜI HƯỚNG DẪN} \\
    \textit{(Ký và ghi rõ họ tên)} \\
    \vspace{2.0cm} \\
    \textbf{PGS. TS. Nguyễn Thái Nghe}
\end{tabular}